\documentclass[11pt]{article}

% --- packages ---------------------------------------------------------------
\usepackage[utf8]{inputenc}
\usepackage[T1]{fontenc}
\usepackage[margin=1in]{geometry}
\usepackage{microtype}
\usepackage{amsmath}
\usepackage{booktabs}
\usepackage{graphicx}
\usepackage{xcolor}
\usepackage{listings}
\usepackage[numbers,sort&compress]{natbib}
\usepackage[colorlinks=true,allcolors=blue]{hyperref}

% --- listings (JSON / ascii diagrams) ---------------------------------------
\definecolor{codebg}{gray}{0.96}
\lstset{
  basicstyle=\ttfamily\footnotesize,
  backgroundcolor=\color{codebg},
  frame=single,
  framesep=4pt,
  rulecolor=\color{gray!40},
  breaklines=true,
  columns=fullflexible,
  keepspaces=true,
  showstringspaces=false,
  aboveskip=1em,
  belowskip=1em,
}

\newcommand{\code}[1]{\texttt{\small #1}}

% Include a screenshot if present, else a labeled placeholder box.
\newcommand{\viewfig}[1]{%
  \IfFileExists{#1}{\includegraphics[width=\linewidth]{#1}}{%
    \fbox{\parbox[c][3cm][c]{0.92\linewidth}{\centering\footnotesize\itshape
      screenshot pending\\\texttt{\scriptsize #1}}}}}

% --- title ------------------------------------------------------------------
\title{\textbf{The Workspace Is the Self-Improving Agent}}

\author{
  Satoshi Nakajima\\
  \small The Singularity Society\\
  \small \texttt{satoshi.nakajima@gmail.com}
}

\date{\today}

\begin{document}
\maketitle
\let\thefootnote\relax\footnotetext{\copyright\ 2026 Satoshi Nakajima. This work is
licensed under a Creative Commons Attribution 4.0 International (CC BY 4.0) license:
\url{https://creativecommons.org/licenses/by/4.0/}.}

\begin{abstract}
A common framing for file-backed LLM assistants is \emph{the workspace is the
database}: the file system stores data and the model is an interface that reads and
writes it. We argue this framing understates what is actually on disk. In
MulmoClaude, an open-source agent application, the workspace holds not only
\emph{data} (notes, records, generated artifacts) but \emph{code}: persona
definitions, tool wiring, application schemas, procedural skills, and recurring
automations, almost all written in domain-specific languages the model interprets.
Once the artifact the agent maintains includes its own \emph{behavior}, the sharper
statement is \textbf{the workspace is the agent}, and an agent that authors its own
workspace is \emph{self-improving}---not because its weights change, but because it
accretes and refines a personal suite of applications in its own substrate. We make
this concrete with three contributions: (i) an account of the workspace-as-agent
architecture, in which identity, capability, behavior, and memory are all
inspectable, forkable (\code{cp -r}), and versionable (\code{git}) files; (ii) a
\emph{self-improvement loop}---recognize, crystallize, tune, retire---by which the
agent promotes a recurring need into a durable application; and (iii) an economic
reframing, \emph{software for an audience of one}, in which the marginal cost of
bespoke software collapses, code becomes a disposable intermediate representation,
and the burden of correctness relocates from a vendor's QA to the trust between one
user and one model. We report a longitudinal single-user deployment in which a
personalized application suite (24 schema-defined applications, 3 external feeds,
and 9 bespoke views) accreted over roughly two months, and we are explicit about
the limitations: self-improvement requires self-pruning, execution is not
replayable, a self-modifying file-based agent widens its own attack surface, and
our evidence is a single author-operated deployment.
\end{abstract}

\noindent\textbf{Keywords:} LLM agents; self-improving agents; agent memory;
end-user software engineering; applications as data; personal software; agent
security.

% ===========================================================================
\section{Introduction}
\label{sec:intro}

A productive design philosophy for file-backed LLM assistants is \emph{the
workspace is the database; files are the source of truth; the model is the
intelligent interface.} It positions the file system as storage and the model as
the component that reads and writes it. The framing is correct but incomplete,
because it describes only half of what lives in a mature agent workspace.

Consider what is actually on disk in MulmoClaude, an open-source agent application
built on the Claude Agent SDK. Alongside the obvious \emph{data}---notes, records,
generated artifacts---are persona definitions that fix the agent's behavior, a
tool-and-server manifest that fixes its capabilities, application schemas that
define entire data-backed apps, procedural skill files the agent reads to operate
those apps, and automation definitions that fix when it acts proactively. These are
not data; they are \emph{program}, almost all written in domain-specific languages
(DSLs) that the model interprets at run time. The model is therefore not merely an
interface onto stored data; it is the \textbf{runtime that interprets a multi-DSL
program persisted as files.} Strip the running process away and the agent does not
vanish---it sits, complete and inspectable, on disk.

That observation motivates a sharper claim: \textbf{the workspace is the agent.}
And because the artifact the agent maintains now includes its own behavior and
capabilities, an agent that edits its workspace is editing itself. We call this
\emph{self-improvement}, and we mean something precise by it (\S\ref{sec:loop}): not
that the model's weights change---they are fixed---but that the agent accretes and
refines a \emph{personal suite of applications} in its own substrate, becoming more
capable in a direction shaped by one particular user.

This paper develops the architecture, the loop that drives it, and its economic
consequence. The detailed mechanism of the application layer---the Collection
DSL, its validator, and the host/model reliability split---is the subject of a
companion paper~\citep{nakajima2026dsl}; here we treat a collection as a given
building block and focus on what it means when the whole agent is files and the
agent authors them.

\paragraph{Contributions.}
\begin{enumerate}
  \item \textbf{Workspace-as-agent.} We show that an agent's identity, capabilities,
  behavior, and memory can all be plain files, making the agent inspectable,
  forkable (\code{cp -r}), and versionable (\code{git}) for free
  (\S\ref{sec:workspace}, \S\ref{sec:forking}).
  \item \textbf{The self-improvement loop.} We describe a recognize $\rightarrow$
  crystallize $\rightarrow$ tune $\rightarrow$ retire loop that promotes a recurring
  need into a durable, self-authored application, yielding personalization through
  accreted applications (\S\ref{sec:loop}).
  \item \textbf{Software for an audience of one.} We argue the marginal cost of
  bespoke software collapses in this setting, that code becomes a disposable
  intermediate representation, and that the cost of correctness relocates rather
  than disappears (\S\ref{sec:one}).
  \item \textbf{A longitudinal case study} of a single real workspace
  (\S\ref{sec:casestudy}), and \textbf{an honest account of the limitations}:
  self-pruning, non-replayability, the security of a self-modifying agent, and the
  single-user scope of our evidence (\S\ref{sec:limitations}).
\end{enumerate}

% ===========================================================================
\section{Background and Related Work}
\label{sec:related}

\paragraph{Self-improving and memory-augmented agents.}
Several systems let an agent extend itself over time. Voyager~\citep{wang2023voyager}
accretes a reusable skill library of generated code; Generative
Agents~\citep{park2023generative} maintain a memory stream they periodically reflect
over; Reflexion~\citep{shinn2023reflexion} and Self-Refine~\citep{madaan2023selfrefine}
improve a task trajectory through verbal self-feedback; STaR~\citep{zelikman2022star}
bootstraps reasoning by fine-tuning on self-generated rationales. These differ from
our setting along two axes. First, \emph{substrate}: the accreted artifact is, variously,
learned weights, a code skill-library, or a memory log, whereas ours is a set of
\emph{declarative DSL documents}---schemas, roles, automations---that are
inspectable, diff-able, and validated before they run. Second, \emph{object of
improvement}: prior work largely improves task-solving \emph{performance}, whereas
the loop here improves the agent's \emph{application surface}---the durable apps it
offers a specific user.

\paragraph{Autonomous agent loops and frameworks.}
AutoGPT~\citep{autogpt} and BabyAGI~\citep{babyagi} popularized autonomous
plan-act-reflect loops; multi-agent frameworks such as AutoGen~\citep{wu2023autogen}
structure collaborating agents with configuration files. These accrete \emph{task
plans} or wire \emph{agent topologies}; our concern is the accretion of durable,
user-facing \emph{applications} into the agent's own substrate, which persist and
are reused long after any single task.

\paragraph{Self-maintained knowledge.}
The direct ancestor of this design is the idea of letting the model build and
maintain its own knowledge base---a growing, interlinked wiki---rather than a flat
memory file~\citep{karpathy2025kb}. MulmoClaude ships exactly such a self-maintained
wiki. This paper extends the principle from self-maintained \emph{data} to
self-maintained \emph{code}: the same workspace also holds the schemas, skills, and
roles the agent authors and refines.

\paragraph{Event-log-primary agents.}
A complementary architecture makes an append-only event log the primary artifact to
obtain deterministic replay and cheap forking~\citep{ynakajima2026log}. We
deliberately choose a different primary artifact---the workspace of files---trading
byte-for-byte run replay (the model step is nondeterministic) for an agent that can
be read, forked, and evolved as plain configuration; \S\ref{sec:limitations} returns
to this trade.

\paragraph{End-user software engineering and the audience of one.}
End-user programming has long sought to let non-programmers build software, from the
spreadsheet onward~\citep{nardi1993small,ko2011enduser}; the spreadsheet was, in
effect, the first widely used \emph{software for an audience of one}. The persistent
price in that tradition is that the user must work inside a constrained authoring
surface and perform the assembly. The shift we describe (\S\ref{sec:one}) keeps the
constrained, validated surface but moves authorship to the agent, and connects this
to a software-economics argument about the collapse of bespoke build cost.

% ===========================================================================
\section{The Workspace Is the Agent}
\label{sec:workspace}

The claim is literal. Table~\ref{tab:contents} inventories a MulmoClaude workspace
by whether each location holds \emph{data} or \emph{code}. Only the bottom rows are
``the database.'' Everything above is program: it specifies what the agent is, what
it can do, how it behaves, and when it acts.

\begin{table}[t]
\centering
\small
\begin{tabular}{@{}lll@{}}
\toprule
\textbf{Location} & \textbf{What it is} & \textbf{Sense} \\
\midrule
\code{config/roles/}          & persona / system-prompt definitions   & code (behavior DSL) \\
\code{config/mcp.json}        & tool-and-server wiring                 & code (capability DSL) \\
collection \code{schema.json} & data model + UI + computation + workflow & code (application DSL) \\
collection \code{SKILL.md}    & how to operate a collection            & code (procedural DSL) \\
\code{data/scheduler/}        & recurring automations                  & code (trigger DSL) \\
\code{data/wiki/}             & densely cross-linked knowledge         & data + graph \\
\code{conversations/memory.md}& accreted facts                         & data \\
\code{artifacts/}             & generated outputs                      & data \\
\bottomrule
\end{tabular}
\caption{What is actually in an agent workspace. The upper rows are program---the
agent's behavior, capabilities, and applications---persisted as files in DSLs the
model interprets. The model is the runtime over this multi-DSL program.}
\label{tab:contents}
\end{table}

The agent's identity is \code{config/roles/}; its capability set is
\code{config/mcp.json} plus its skills; its applications are its collections; its
proactive behavior is its automations; its long-term knowledge is its wiki and
memory. None of this is a learned blob separate from the files: the agent
\emph{is} the files, of which data is one part.

\paragraph{Lineage: from self-maintained memory to self-improvement.}
The design descends from the self-maintained knowledge base~\citep{karpathy2025kb}:
let the model curate its own artifact rather than be handed one. A wiki is
self-maintained \emph{data}; a collection extends that to self-maintained
\emph{structured} data; the observation here is that the workspace also holds
self-maintained \emph{code}---schemas, skills, roles, automations. The moment the
maintained artifact includes the agent's own behavior, ``self-maintaining'' becomes
``self-improving'' (\S\ref{sec:loop}). This kind of improvement \emph{compounds}:
each application the agent crystallizes makes the next interaction faster and more
reliable, and the workspace grows into a personal capability surface no general
model update reproduces.

% ===========================================================================
\section{A Collection Is an Application}
\label{sec:collection}

The load-bearing building block is the \emph{collection}: a single
\code{schema.json} plus record files that, together with the model as runtime,
behaves like a traditional application. It has a data model (typed fields),
relationships (cross-collection references), computation (spreadsheet-like derived
fields), a user interface (host-rendered field types plus optional bespoke HTML
views), workflows (declarative actions that hand off to a seeded chat), and
host-driven recurrence---with \emph{zero} application-specific host code. The
companion paper~\citep{nakajima2026dsl} develops this mechanism and the validator
that makes a structurally incoherent schema fail to load; for this paper the
relevant point is only that a collection is a genuine application, not a memory
note. It differs from a note in three ways: \emph{structure} (a schema pins the
shape; a note is re-parsed lossily), \emph{durable behavior} (a collection behaves
the same tomorrow; a prompt-of-the-moment does not), and \emph{a UI for free}. That
makes a collection the top rung of a ladder---conversation, then \code{memory.md},
then a collection, then a collection with workflows and UI---and recognizing when a
recurring need has stabilized enough to climb it is where self-improvement begins.

% ===========================================================================
\section{The Boundaries That Dissolve}
\label{sec:boundaries}

A traditional application is built by developers, shipped across a deployment
boundary, and run by users. The workspace model erases all three.

\paragraph{The builder is the agent, not a human composer.}
This is the line that separates the architecture from no-code platforms. In
Airtable, Notion, or Retool---even with recent AI-assisted authoring---the human
still composes the app from primitives the platform supplies. Here the user
expresses a need in natural language, often implicitly through repeated similar
requests, and the \emph{agent} writes the schema, authors the skill, and wires the
actions. The human is the product owner, not the carpenter: the jump is from
\emph{no-code} (a human composes blocks) to \emph{intent-driven} (an agent composes
the DSL program from intent).

\paragraph{There is no deployment boundary.}
A collection is authored directly into the agent's own substrate; building it and
running it happen in the same workspace, in the same breath. There is no build, no
CI, no release, no migration job. ``Also track whether I have been to each
restaurant'' edits the schema and skill in place, and the app has evolved, live.

\paragraph{The app/agent boundary disappears.}
Because the workspace is the agent (\S\ref{sec:workspace}) and the app is authored
into the workspace, the app becomes part of the agent. There is no longer a clean
line between ``the agent'' and ``the applications the agent built.'' When the agent
refines a collection it built, it is modifying its own code. That is the precise
sense in which this is a \emph{self-improving agent}, and why that description fits
better than ``no-code platform'': the latter foregrounds a human builder and a
separate product, both of which are wrong here.

\paragraph{The trajectory already shows it.}
This is not aspirational. In the studied codebase, features such as worklog, client,
and invoice tracking began as hardcoded host plugins---developer-built apps behind a
deployment boundary---and are being superseded by collection skills the agent can
author with no host change. The locus of application creation is visibly moving from
the development team to the agent.

% ===========================================================================
\section{The Self-Improvement Loop}
\label{sec:loop}

``Self-improving'' is loose unless we say what improves and how. It is not that the
model gets generally smarter; it is that the agent accretes and refines a personal
suite of applications. The loop has four steps, and the fourth is the one most
easily forgotten.

\begin{lstlisting}
   RECOGNIZE          CRYSTALLIZE         TUNE              RETIRE
   a recurring   -->  into a DSL    -->   the host/model -->  when
   need               application         split (dial)        superseded
     ^                                                          |
     +----------------------------------------------------------+
\end{lstlisting}

\begin{enumerate}
  \item \textbf{Recognize.} Notice that a need recurs---the user keeps asking for the
  same shape of thing in free-form chat. This is the signal to promote it up the
  maturity ladder (\S\ref{sec:collection}).
  \item \textbf{Crystallize.} Author the application: write the schema, the skill,
  the actions. The loose intent becomes a durable, structured, reliable app.
  \item \textbf{Tune.} As the need evolves, adjust where the reliability dial sits
  (companion paper~\citep{nakajima2026dsl})---move more structure into the schema
  where it must be deterministic, leave more to the model where it needs judgment.
  \item \textbf{Retire.} When an app is superseded or abandoned, remove it. This
  keeps the agent's ``self'' coherent rather than sprawling (\S\ref{sec:limitations}).
\end{enumerate}

Running this loop over months yields not a smarter model but an agent whose
workspace has grown into a \emph{personalized application suite shaped to one user}.
Two instances that start identical diverge into very different agents, because they
crystallize different apps from different lives. Self-improvement here is
\textbf{personalization through accreted applications}---the agent becomes more
useful to \emph{its} owner specifically, not more capable in the abstract.

% ===========================================================================
\section{Software for an Audience of One}
\label{sec:one}

The self-improvement loop is only affordable because of an economic shift worth
stating on its own terms.

\paragraph{The collapse of bespoke build cost.}
Software was always built for the average of an audience, never the individual---but
this was a consequence of cost, not a preference. Development was expensive and
fixed while replication was free, so the rational strategy was to spread the fixed
cost over as many users as possible: build for the mass. When an agent can generate
a working application from a sentence of intent, the fixed cost collapses, and with
it the reason to target the average. Software whose audience is a single
person---historically impossible, not merely uneconomic---becomes the default for a
large class of needs.

\paragraph{Code as a disposable intermediate representation.}
When code costs almost nothing to produce, it stops being the precious asset and
becomes a compile target between intent and behavior. The durable artifacts are the
\emph{intent} (natural language) and the \emph{data} (the records); the code in
between---a schema, or the HTML of a custom view---is regrown when the need changes
rather than maintained. Figure~\ref{fig:views} shows three custom views from the
studied workspace: each is a self-contained HTML file the agent wrote on request and
the host serves sandboxed, and in each case the owner described the view in natural
language and never saw the markup. A restaurant guide drawn as a transit map is
software no vendor would build for a market of one---and now does not have to be.

\begin{figure}[t]
\centering
\begin{minipage}[t]{0.31\textwidth}\centering
  \viewfig{figures/restaurants-map.png}\\[2pt]\footnotesize (a) a restaurant guide as a transit-style map
\end{minipage}\hfill
\begin{minipage}[t]{0.31\textwidth}\centering
  \viewfig{figures/portfolio-allocation.png}\\[2pt]\footnotesize (b) a portfolio as an allocation chart
\end{minipage}\hfill
\begin{minipage}[t]{0.31\textwidth}\centering
  \viewfig{figures/feed-gallery.png}\\[2pt]\footnotesize (c) a podcast feed as a read--write gallery
\end{minipage}
\caption{Three agent-authored custom views from the studied deployment (two
collection views and one feed view). Each is software with an addressable audience
of one, produced for the cost of a natural-language request; the owner never wrote
or read the generated HTML.}
\label{fig:views}
\end{figure}

\paragraph{Correctness relocates rather than disappears.}
Mass software is poor at fitting the individual but excellent at accumulating
correctness: many users and years surface edge cases and harden the code. Software
for an audience of one has no such second audience. The cost of correctness does not
vanish; it relocates---from a vendor's QA organization to the trust between one user
and one model. The stakes are uneven: a wrong poster gallery is harmless, but a
wrong financial or tax view is not, and only the model and the single user ever
inspect it. This is the spreadsheet era's oldest hazard---the invisible formula error
in a model no one else audits---returning in a more powerful form, and it is a real
limitation (\S\ref{sec:limitations}), not a footnote.

% ===========================================================================
\section{Case Study: A Longitudinal Single-User Deployment}
\label{sec:casestudy}

We report on one author's production workspace, which accreted applications through
ordinary use over roughly two months, with no host code written per application.

\paragraph{An accreted, personal application suite.}
The workspace contains 24 schema-defined collections holding $\sim$700 records,
3 external feeds holding a further $\sim$260 ingested records, and 9 agent-authored
custom views (7 over collections, 2 over feeds). The applications span personal
finance, reading and watch lists, a vocabulary trainer, recurring-obligation
trackers, contact and engagement records, and domain-specific trackers; their record
counts span nearly three orders of magnitude, from a single-record profile to
hundreds of entries. Crucially, the suite \emph{diverges sharply from any
mass-market product}---a restaurant guide rendered as a transit map, a read--write
podcast tracker carrying the owner's own ratings and resume positions---which is the
signature of build-for-one (\S\ref{sec:one}).

\paragraph{The loop, observed.}
All four steps of \S\ref{sec:loop} appear in the history. Needs were
\emph{recognized} and \emph{crystallized} into collections; schemas were
\emph{tuned} in place (for example, a podcast feed gained user rating and notes
fields after ingestion had populated it, absorbed as an additive, non-breaking
change); and the \emph{retire} step is visible as the ongoing migration from
hardcoded worklog/client/invoice plugins toward collection skills that supersede
them (\S\ref{sec:boundaries}).

\paragraph{Forking and versioning, in practice.}
Because the agent is its files, the whole agent---apps, memory, personality,
automations---is copied by \code{cp -r} and versioned by \code{git}; every
self-edit to a schema, role, or skill is a readable, revertable diff. This is the
audit mechanism we rely on in place of run replay (\S\ref{sec:limitations}).

% ===========================================================================
\section{Forking, Portability, and What the Workspace Buys for Free}
\label{sec:forking}

Because the agent is its workspace, several properties that normally require
engineering fall out of the file system. \emph{Forking} an agent is \code{cp -r}:
copying the directory copies the apps, memory, personality, and automations as one
coherent unit, with no separate ``export behavior'' and ``export data'' steps,
because they are the same files. \emph{Versioning} an agent is \code{git}: a
git-backed workspace gives configuration replay (\code{checkout}), branching, and
self-edit lineage (\code{blame})---audit and rollback for the agent's edits to
itself, with no bespoke runtime. \emph{Portability} is trivial: nothing is stranded
in a process or a database server; move the directory, move the agent. And
\emph{inspectability} is total: every part of the agent is a readable, greppable,
diffable file, with no opaque learned blob that constitutes the agent apart from its
workspace. These are not the headline result, but they are why the headline result
is practical: an agent you can fork, version, inspect, and carry is one you can more
safely let modify itself, because every modification is a visible, reversible change
to files you can read.

% ===========================================================================
\section{Discussion and Limitations}
\label{sec:limitations}

\paragraph{Self-improvement requires self-pruning.}
If every recognized need spawns an artifact and every artifact becomes part of the
agent's ``self,'' then self-improvement without self-pruning produces a junk drawer:
half-built collections, abandoned skills, stale automations. Traditional
applications have maintainers and sunset processes; a self-built suite needs the
same discipline, and nothing provides it automatically. The retire step
(\S\ref{sec:loop}) is what keeps the agent's self coherent rather than merely
larger; an agent that only ever adds capabilities eventually degrades, because more
surface means more ways to be wrong and more instructions competing for the same
intent. A genuinely self-improving agent must be as willing to delete its own code
as to write it.

\paragraph{Non-replayable execution.}
Because the operating step is a nondeterministic model call, a run cannot be replayed
byte-for-byte, and a self-modification's effect cannot be A/B-replayed against the
prior behavior on identical inputs---the deliberate trade against log-primary
designs~\citep{ynakajima2026log}. The mitigation is configuration-level rather than
run-level: a git-backed workspace makes every self-edit auditable and revertable,
which is the granularity that matters for an agent editing its own applications.

\paragraph{Correctness has no second audience.}
As argued in \S\ref{sec:one}, an application with an audience of one does not
accumulate correctness the way mass software does. Pushing consequential computation
into host-evaluated derived fields helps, but relocates rather than removes the
valid-but-wrong failure mode: a derived formula is itself a model-authored
expression that the host evaluates faithfully whether or not it is the right formula.
That formulas and schemas are short, inspectable, plain-text artifacts makes
\emph{review against stated intent} tractable, but the residual risk is real.

\paragraph{The security of a self-modifying, file-based agent.}
Making behavior editable as files widens the attack surface in two ways. First,
content the agent ingests is untrusted: feeds bring external natural language into
records the model reasons over while it may hold write-capable tokens, so indirect
prompt injection~\citep{greshake2023injection} is a live threat that view sandboxing
does not address, since it governs UI rendering rather than the model's context.
Second, the agent edits its \emph{own} behavior---roles, skills, schemas---so a
successful injection or a careless self-edit can in principle alter what the agent is,
not just what it stores. The defenses are partial and must be layered: a
git-backed workspace makes every self-edit reviewable and revertable; the schema
validator rejects structurally incoherent self-modifications at load regardless of
how they were written; and consequential or externally-triggered operations warrant
confirmation. We did not observe an attack in our benign deployment, but we flag
self-modification security as the architecture's most important open problem.

\paragraph{Threats to validity.}
Our evidence is a single user, a single system, and an author-operated deployment;
it demonstrates feasibility and the no-host-code property but does not establish
generality, and it reports usage rather than a controlled comparison. Quantifying
how often the loop is run, how much authoring effort it saves versus alternatives,
and how the discipline of pruning holds up across users and longer horizons is
future work.

\paragraph{Owning the learning loop.}
A firm's strategic question in an AI-driven economy---recently framed by
Nadella~\citep{nadella2026ecosystem}---is whether it can \emph{own the learning
loop} that encodes its institutional knowledge, and, as the test of that ownership,
whether it can swap out a generalist model without losing the ``company-veteran''
expertise built on top of it. The workspace-as-agent architecture passes that test
by construction: the expertise lives in roles, schemas, skills, and the
wiki---files the model interprets---not in weights it carries, so replacing the
model leaves the accumulated knowledge intact. This is a stronger form of ownership
than encoding the loop by fine-tuning a model on private traces, which re-couples
the firm's IP to a particular model; a declarative workspace keeps that IP
inspectable, diff-able, forkable, and model-agnostic. The honest qualification is
one of coverage: declarative files capture \emph{explicit} knowledge---workflows,
judgment-as-prose, linked notes---better than the \emph{tacit} pattern recognition
that learning on private traces can absorb. The two are complementary; the workspace
is the portable substrate that survives a model swap, and any private tuning is an
accelerator layered on top of it rather than the thing one owns.

\paragraph{From one user to the firm.}
MulmoClaude is, throughout this paper, a \emph{single-user} agent, and our evidence
is one personal workspace---but nothing in the architecture is intrinsically
personal, and the concept generalizes directly to an organization. The same
primitives scale: a team's learning loop becomes a shared, versioned workspace of
roles, collection schemas, skills, and a cross-linked wiki, forked with \code{cp -r},
branched and audited with \code{git}, and owned independently of any model vendor.
At that scale the workspace is exactly the firm's compounding capital---institutional
memory made queryable, recurring workflows crystallized into applications, and the
whole of it portable across model generations. We have not evaluated the multi-user
case, and its open problems are real (concurrency, access control, and the merge
discipline that keeps a shared ``self'' coherent); but the single-user results
suggest that the durable unit of ownership---whether the owner is a person or a
firm---can be a workspace of files rather than a tuned model.

% ===========================================================================
\section{Conclusion}
\label{sec:conclusion}

Traditional software assumes that engineers write programs and users operate them.
The architecture described here assumes something different: users express intent,
the agent writes the programs---into itself---and the resulting applications are part
of the agent. The collection schema is the program, the DSLs in the workspace are
the source code, the model is the runtime, and the workspace is the agent. Because
the programs live inside the agent, writing them is the agent improving itself, which
is why, of the available names, \emph{self-improving agent} fits better than
\emph{no-code platform}. The same move that makes this possible---near-zero-cost
bespoke software, authored from intent---turns software from a product shipped to a
market into an artifact grown for an individual, and shifts the discipline from
maintaining precious code to verifying disposable code and curating the durable
intent and data behind it.

% ===========================================================================
\section*{Availability}
MulmoClaude is open source, including the roles, collection schemas, skills, and
automations that constitute the workspace-as-agent described here, and the
self-maintained wiki: \url{https://github.com/receptron/mulmoclaude}.

\bibliographystyle{plainnat}
\bibliography{refs}

\end{document}
